
% =========================================================
% SLIDE 2 - BACKGROUND & PROBLEM STATEMENT
% =========================================================
\begin{frame}{Bối cảnh và vấn đề nghiên cứu}

\small

\textbf{\color{HCMUTBlue}Bối cảnh (Context)}

\vspace{0.15cm}

Trong môi trường đại học, sinh viên cần sử dụng các
phần mềm chuyên dụng như \textbf{MATLAB, SolidWorks,
Oracle DB, SQL Server} để học tập và thực hành.


Tuy nhiên, nhu cầu sử dụng không diễn ra đồng thời
và liên tục. Khiến việc đầu tư tài nguyên và license
theo tổng số sinh viên có thể gây lãng phí.

\vspace{0.35cm}

\textbf{\color{HCMUTBlue}Giải pháp đề xuất (Proposed Solution)}

\vspace{0.15cm}

Một giải pháp tiềm năng là xây dựng
\textbf{phòng lab ảo trên nền tảng Cloud},
trong đó sinh viên được cấp môi trường thực hành
thông qua \textbf{Virtual Machine (VM)} theo nhu cầu.

\vspace{0.35cm}

\textbf{\color{HCMUTBlue}Bài toán đặt ra (Problem Statement)}

\vspace{0.1cm}

\begin{itemize}
  \setlength{\itemsep}{0.45em}

  \item Duy trì nhiều VM liên tục
  $\rightarrow$ lãng phí tài nguyên và chi phí.

  \item Chỉ khởi động VM khi có yêu cầu
  $\rightarrow$ phát sinh thời gian chờ.

  \item Nhu cầu tăng đột biến
  $\rightarrow$ nguy cơ quá tải và hàng chờ dài.

  \item Tài nguyên và license có giới hạn
  $\rightarrow$ cần cơ chế phân bổ hợp lý.
\end{itemize}

\vfill

\begin{center}
  \color{HCMUTBlue}
  \textbf{Mục tiêu:}
  Cân bằng giữa \textbf{chi phí tài nguyên}
  và \textbf{thời gian chờ của sinh viên}.
\end{center}

\end{frame}
